\subsection{Simple Regression Using Frequency}

Simple linear regression was performed using both Python's
\texttt{statsmodels} package and ScalaTion. Both implementations produced
essentially identical coefficient estimates. The fitted model was

\[
\widehat{\mathrm{SPL}}
=
127.3037 - 0.000855(\mathrm{Frequency}).
\]

The negative slope indicates an inverse association between frequency and
sound pressure level. An increase of 1,000 Hz is associated with an estimated
decrease of approximately 0.855 dB in sound pressure level. The slope was
statistically significant ($p<0.001$).

ScalaTion reported $R^2=0.152655$ and adjusted $R^2=0.152091$, which agreed
with the corresponding statsmodels results to rounding precision. Thus,
frequency alone explains approximately 15.3\% of the observed variability in
sound pressure level. Although the relationship is statistically significant,
most of the variation in the response remains unexplained by frequency alone.

\subsection{Simple Regression Using Displacement Thickness}

Simple linear regression was also performed using displacement thickness as
the predictor. Once again, ScalaTion and statsmodels produced nearly identical
coefficient estimates. The fitted equation was

\[
\widehat{\mathrm{SPL}}
=
126.6632 - 164.0275(\mathrm{Displacement\ Thickness}),
\]

where displacement thickness is measured in meters.

The fitted slope of $-164.0275$ dB per meter is equivalent to approximately
$-0.164$ dB per millimeter. Therefore, a 1 mm increase in displacement
thickness is associated with an estimated decrease of approximately 0.164 dB
in sound pressure level. The coefficient was statistically significant
($p<0.001$).

ScalaTion reported $R^2=0.097762$ and adjusted $R^2=0.097161$, again matching
the statsmodels results to rounding precision. Thus, displacement thickness
alone explains approximately 9.8\% of the observed variability in sound
pressure level.

\subsection{Regression Discussion}

Both selected predictors had statistically significant negative relationships
with sound pressure level, and both software packages produced essentially
identical results. Frequency provided the stronger simple regression model,
with an $R^2$ of approximately 0.153 compared with approximately 0.098 for
displacement thickness. Nevertheless, neither predictor explains a large
proportion of the response variability by itself. These results suggest that
sound pressure level is influenced by additional experimental variables and
that a model incorporating multiple predictors may provide substantially
greater explanatory power.
